\section{Compliance with Standards and Professional Practice}

The project follows software engineering and professional practice principles relevant to educational AI systems. It emphasizes honest reporting of results, clear privacy-scope limitations, responsible handling of protected assessment material, and reproducible evaluation artifacts.

\subsection{Software Standard}
The system follows modular software design. Ingestion, retrieval, privacy screening, generation, evaluation, and user-interface responsibilities are separated. Results are stored in structured JSON/CSV files for auditability. The report cites external research and avoids unsupported claims.

\subsection{Data Privacy and Security Practice}
The project handles protected exam artifacts as confidential teacher-only resources. Student-facing retrieval is restricted to public material. The privacy guard checks user queries and outputs for reconstruction risk. The system does not claim complete security, differential privacy, or protection against compromised accounts.

\subsection{Professional Ethics}
The project is designed to support learning without compromising exam fairness. It respects intellectual property by citing external work and avoids presenting generated answers as authoritative when retrieved evidence is insufficient.

\section{Design Constraints}

\subsection{Ethical and Professional Responsibilities}
The main ethical responsibility is to prevent academic-assessment leakage while still supporting student learning. The project addresses this through role-separated corpora, protected metadata, output screening, and conservative claim reporting.

\subsection{Environmental and Sustainability Considerations}
The project uses a quantized local model below 1 GB and CPU-oriented inference. This reduces reliance on large cloud inference for small course-support tasks. However, larger deployments should monitor energy consumption and server utilization.

\subsection{Health and Safety Considerations}
The system does not control physical devices, so direct physical safety risk is low. The main safety issue is information safety: protected exam content must not be exposed to unauthorized users. The role-aware workflow and privacy guard mitigate this risk.

\subsection{Economic Constraints}
The system is designed to use open-source tools and existing hardware. This reduces deployment cost for universities with limited infrastructure.

\section{Complex Engineering Problem (CEP) Coverage}

The project qualifies as a complex engineering problem because it involves conflicting requirements, multiple interacting modules, privacy-sensitive academic data, domain shift, and measurable trade-offs between utility and protection.

\begin{table}[H]
\centering
\caption{Mapping with Complex Engineering Problems (Ps)}
\label{tab:cep_mapping}
\begin{tabular}{|p{0.08\textwidth}|p{0.22\textwidth}|p{0.43\textwidth}|p{0.19\textwidth}|}
\hline
\textbf{Ps} & \textbf{Attributes} & \textbf{How Ps are addressed through the project} & \textbf{Evidence}\\
\hline
P1 & Depth of knowledge required & The project combines RAG, vector retrieval, privacy screening, Bloom classification, OCR, and local LLM inference. & Chapters 2--5\\
\hline
P2 & Conflicting requirements & Student utility conflicts with protected exam confidentiality. The design balances benign allow rate and attack block rate. & Sections 3.1, 5.3\\
\hline
P3 & Depth of analysis & The system is evaluated using Bloom metrics, privacy prompt taxonomy, QA/RAG metrics, and OCR metrics. & Chapter 5\\
\hline
P4 & Familiarity of issues & RAG privacy in the university assessment setting is not a routine single-module problem. & Sections 2.4, 3.3\\
\hline
P5 & Extent of applicable codes & The project follows software ethics, privacy-aware design, and evidence-aligned reporting principles. & Chapter 7\\
\hline
P6 & Stakeholder involvement & Students, teachers, administrators, and system operators have different roles and requirements. & Section 3.1\\
\hline
P7 & Interdependence & Ingestion metadata affects retrieval, privacy screening, generation, and evaluation. & Chapters 3--4\\
\hline
\end{tabular}
\end{table}

\section{Complex Engineering Activities (CEA) Coverage}

The project also involves complex engineering activities because it requires coordination of multiple technical resources, software modules, datasets, and evaluation methods.

\begin{table}[H]
\centering
\caption{Mapping of Complex Engineering Activities (As) to Project Implementation}
\label{tab:activities_mapping}
\begin{tabular}{|p{0.08\textwidth}|p{0.22\textwidth}|p{0.43\textwidth}|p{0.19\textwidth}|}
\hline
\textbf{As} & \textbf{Attributes} & \textbf{How As are addressed through the project} & \textbf{Evidence}\\
\hline
A1 & Range of resources & Uses datasets, OCR backend, FAISS, local GGUF model, evaluation scripts, and LaTeX reporting. & Chapters 4--6\\
\hline
A2 & Interaction complexity & Public/protected corpora, role policies, retrieval, and output screening interact closely. & Chapter 3\\
\hline
A3 & Innovation & Integrates student RAG, teacher moderation, Bloom analysis, and privacy screening in one local workflow. & Chapters 3--5\\
\hline
A4 & Consequences & Failure could expose exam artifacts or block legitimate learning support. & Sections 3.3, 5.3\\
\hline
A5 & Familiarity & The combined university privacy-constrained RAG workflow is not a routine CRUD application. & Chapters 1--2\\
\hline
\end{tabular}
\end{table}

\section{Summary}

This chapter mapped the project to professional practice, design constraints, CEP attributes, and CEA activities. The project demonstrates complex engineering characteristics through its privacy-utility trade-off, multimodal ingestion, local Qwen deployment (LoRA + GGUF), and role-aware academic workflow.
